\begin{table}[htbp]
\centering
\caption{Counterfactual Policy Designs: Exempting High-Value Items}
\label{tab:counterfactual}
\begin{adjustbox}{max width=\textwidth}
\begin{threeparttable}
\small
\begin{tabular}{lccccc}
\toprule
Policy design & Quartiles covered & $\hat\beta$ range & Items kept & Fiscal cost (R\$~M) & Savings (R\$~M) \\
\midrule
\multicolumn{6}{l}{\textit{Panel A: DiD price coefficient by Group-65 reference-value quartile}} \\
Q1 (R\$0.00 -- R\$500.00) & 1 & -0.0967*** & 262,111 & -- & -- \\
Q2 (R\$500.00 -- R\$1875.00) & 2 & -0.1151*** & 172,531 & -- & -- \\
Q3 (R\$1875.00 -- R\$7622.91) & 3 & -0.1224*** & 123,742 & -- & -- \\
Q4 (R\$7622.91 -- $\infty$) & 4 & -0.0967*** & 91,330 & -- & -- \\
\midrule
\multicolumn{6}{l}{\textit{Panel B: Fiscal cost under alternative policy designs}} \\
Current (no exemption) & Q1,2,3,4 & -- & 65,606 & 64.7 & 0.0 \\
Exempt top quartile (Q4) & Q1,2,3 & -- & 49,204 & 6.2 & 58.4 \\
Exempt top half (Q3, Q4) & Q1,2 & -- & 32,798 & 1.5 & 63.2 \\
Keep only bottom quartile (Q1) & Q1 & -- & 16,351 & 0.2 & 64.4 \\
\bottomrule
\end{tabular}
\begin{tablenotes}
\small
\item \textit{Notes:} Panel~A estimates the DiD price coefficient separately
within each quartile of pre-period Group-65 reference value (quartile
cutoffs listed in the first column). Panel~B simulates fiscal cost under
alternative policy designs that exempt high-value items from the SME-only
rule. Fiscal cost is computed as
$\sum_{q \in \text{covered}} |e^{\hat\beta_q} - 1| \times V_q$
where $V_q$ is the Group-65 pre-period procurement value in quartile~$q$.
Total pre-period value: R\$689.0 million across 65,606 items.
Exempting the top quartile alone removes most of the fiscal cost while
retaining the ME/EPP preference for the items where heterogeneity analysis
shows the policy is least distortive.
*** \textit{p}$<$0.01, ** \textit{p}$<$0.05, * \textit{p}$<$0.1.
\end{tablenotes}
\end{threeparttable}
\end{adjustbox}
\end{table}
